% - Preambulo
\documentclass[12pt, letterpaper]{report}

%%% - Paquetes
\usepackage{babel}
\usepackage[utf8]{inputenc}
\usepackage{xcolor}
\usepackage{pifont}
\usepackage{lettrine}
\usepackage{lmodern}
\usepackage{afterpage}
\usepackage{graphicx}
\graphicspath{ {images/} }

%%% - Definimos la zona de la pagina ocupada por el texto
\oddsidemargin -1.0cm
\headsep -1.0cm
\textwidth=18.5cm
\textheight=23cm

%%% - Definicion de color personalizado
\definecolor{gold}{RGB}{255, 219, 87}

%%% - Definmos el tiempo de salto de linea
\setlength{\parskip}{\baselineskip}

%%% - Definimos TÍTULO
\title{
	\Huge {\bfseries \textcolor{gold}{DOMINGO DE RAMOS\\ EN LA PASIÓN DEL SEÑOR}} \\
	\vspace{2em}
	\includegraphics{ancla}
}

\author{
	\large {\bfseries \textcolor{gold}{Conmemoración de la entrada del Señor en Jerusalén}}
}

\date{ }

% - Cuerpo
\begin{document}
%%%%%%%%%%%%%%%%%%%% P R I N C I P I O %%%%%%%%%%%%%%%%%%%%%%%%%%%%%%
	\pagecolor{red}\afterpage{\nopagecolor}
	% Generamos título
	\maketitle

	\begin{center}
	\Huge {\bfseries PROCESIÓN}
	\end{center}

	\large {\textcolor{red}{El Presbítero, revestido con las vestiduras rojas que se requieren para la celebración de la misa, se dirigen al lugar donde se ha congregado el pueblo y estan ya los ramos y todo lo necesario para bendecirlos.}}

	\Large {\bfseries \textcolor{red}{Saludo}}\\
	\large {\textcolor{red}{El sacerdote y el pueblo se signan, mientras el sacerdote dice:}}

	\noindent
	\Large {\bfseries \textcolor{red}{V/.}} \hspace{0.5cm} {En el nombre del Padre y del Hijo \Huge{\textcolor{red}{\ding{64}}} \Large y del Espíritu Santo.}\\
	\Large{{\bfseries \textcolor{red}{R/.}} \hspace{0.5cm} Amén.}

	\noindent
	\Large {{\bfseries \textcolor{red}{V/.}} \hspace{0.5cm} La gracia de nuestro Señor Jesucristo,\\
	. \hspace{1.5cm} el amor del Padre y la comunión del Espíritu Santo\\
	. \hspace{1.5cm} esté con todos vosotros.}\\
	\Large{{\bfseries \textcolor{red}{R/.}} \hspace{0.5cm} Y con tu espíritu.}

	\Large {\bfseries \textcolor{red}{Monición}}\\
	\large {\textcolor{red}{Después hace una breve monición, en la que invita a los fieles a participar activa y conscientemente en la celebración de este día, con estas palabras u otras semejantes:}}

	\Large {Queridos hermanos:\\
	Ya desde el principio de la Cuaresma nos venimos preparando con obras de penitencia y caridad. Hoy, cercana ya la Noche santa de la Pascua, nos disponemos a inaugurar, en comunión con toda la Iglesia, la celebración anual de los misterios de la pasión y resurrección de Jesucristo, misterios que empezaron con la solemne entrada de Jesús en Jerusalén.}

	\Large {Por ello, recordando con fe y devoción la entrada triunfal de Jesucristo en la ciudad santa, le acompañaremos con nuestros cantos, para que, participando ahora de su cruz, merezcamos un día tener parte en su resurrección.}

	\newpage

	\Large {\bfseries \textcolor{red}{Bendición de los ramos}}\\
	\large {\textcolor{red}{Después de la monición, el Presbítero bendice los ramos con la siguiente plegaria:}}

	\noindent
	\lettrine[lines=2]{\bfseries \textcolor{red}{D}}{} \Large {ios todopoderoso y eterno,\\
	santifica con tu} \Huge{\textcolor{red}{\ding{64}}}  \Large{bendición estos ramos,\\
	y, a cuantos vamos a acompañar a Cristo Rey aclamándolo con cantos,\\
	concédenos, por medio de él, entrar en la Jerusalén del cielo.\\
	Él, que vive y reina por los siglos de los siglos.}

	\noindent
	\Large {\bfseries \textcolor{red}{R/.}} \hspace{1cm} Am\'en

	\large {\textcolor{red}{A continuación se acerca a los ramos y los asperja con agua bendita sin decir nada.}}

	\large {\textcolor{red}{Seguidamente, el Presbítero proclama el Evangelio. Puede utilizarse incienso, si se juzga oportuno. Después del evangelio, se puede hacer una breve homilía.}}

	\large {\textcolor{red}{Antes de comenzar la procesión, el Presbítero, el diácono o un ministro laico, dice con estas u otras palabras:}}

	\Large {Queridos hermanos:\\
	Imitando a la multitud que aclamaba al Señor, \\
	acompañemoslo también nosotros con júbilo.}

	\large {\textcolor{red}{Y comienza la procesión hacia la iglesia donde se va a celebrar la misa.}}

	\newpage

	\begin{center}
	\Huge {\bfseries MISA}
	\end{center}

	\large {\textcolor{red}{Después de la procesión el Presbítero comienza la misa con la oración colecta.}}

	\Large {\bfseries \textcolor{red}{Oraci\'on colecta}}
	\lettrine[lines=2]{\bfseries \textcolor{red}{D}}{}\Large {ios todopoderoso y eterno,\\
	que hiciste que nuestro Salvador se encarnase y soportara la cruz\\
	para que imitemos su ejemplo de humildad, concédenos, propicio,\\
	aprender las enseñanzas de la pasión y participar de la resurrección gloriosa.\\
	Por nuestro Señor Jesucristo.}

	\noindent
	\Large{{\bfseries \textcolor{red}{R/.}} \hspace{0.5cm} Amén.}

	\begin{center}
	\Huge {\bfseries Lit\'urgia de la Palabra}
	\end{center}

	\large {\textcolor{red}{Para la lectura de la Pasión del Señor no se llevan ni cirios ni incienso, ni se hace al principio la salutación habitual, ni se signa el libro. Esta lectura la proclama un diácono o, en su defecto, el mismo celebrante. Puede también ser proclamada por lectores laicos, reservando, si es posible, al sacerdote la parte correspondiente a Cristo. Si los lectores de la Pasión son diáconos, piden, como de costumbre, la bendición del celebrante antes de empezar la lectura; pero si los lectores no son diáconos se omite esta bendición.}}

	\Large {\bfseries \textcolor{red}{Homil\'ia}}\\
	\large {\textcolor{red}{Después de la lectura de la historia de la Pasión téngase, oportunamente, una breve homilía. También puede observarse algún espacio de silencio.}}

	\Large {\bfseries \textcolor{red}{Credo}}\\
	\large {\textcolor{red}{Acabada la homil\'ia, si la liturgia del d\'ia lo prescribe, se hace la profesi\'on de fe:}}

	\noindent
	\Large {Creo en un solo Dios...}

	\Large {\bfseries \textcolor{red}{Oraci\'on de los fieles}}\\
	\large {\textcolor{red}{Seguidamente se tiene la oraci\'on de los fieles.}}

	\begin{center}
		\Huge {\bfseries Lit\'ugia Eucar\'istica}
	\end{center}

	\Large {\bfseries \textcolor{red}{Presentación de Ofrendas}}\\
	\large {\textcolor{red}{Las ofrendas son llevadas de la crencia al altar. Despu\'es, de pie en el centro del altar extendiendo y juntando las manos, dice:}}

	\lettrine[lines=2]{\bfseries \textcolor{red}{O}}{} \Large ren, hermanos, para que este sacrificio, m\'io y de ustedes,\\
	sea agradable a Dios, Padre todopoderoso.

	\noindent
	\Large {\bfseries \textcolor{red}{R/.}} \hspace{0.5cm} El Se\~nor reciba de tus manos este sacrificio para alabanza y gloria de su nombre, para nuestro bien y el de toda su santa Iglesia.

	\Large {\bfseries \textcolor{red}{Oraci\'on sobre las ofrendas}}
	\lettrine[lines=2]{\bfseries \textcolor{red}{S}}{}\Large eñor, que por la pasión de tu Unigénito\\
	se extienda sobre nosotros tu misericordia y,\\
	aunque no la merecen nuestras obras,\\
	que con la ayuda de tu compasión\\
	podamos recibirla en este sacrificio único.\\
	Por Jesucristo, nuestro Señor.

	\noindent
	\Large{{\bfseries \textcolor{red}{R/.}} \hspace{0.5cm} Amén.}

	\begin{center}
		\Large PLEGARIA EUCARISTICA II
	\end{center}

	\Large {\bfseries \textcolor{red}{PREFACIO de la Pasión del Señor}}

	\noindent
	\Large {\bfseries \textcolor{red}{V/.}} \hspace{0.5cm} El Se\~nor est\'e con vosotros.\\
	\noindent
	\Large {\bfseries \textcolor{red}{R/.}} \hspace{0.5cm} Y con tu esp\'iritu.

	\noindent
	\Large {\bfseries \textcolor{red}{V/.}} \hspace{0.5cm} Levantemos el coraz\'on.\\
	\noindent
	\Large {\bfseries \textcolor{red}{R/.}} \hspace{0.5cm} Lo tenemos levantado hacia el Se\~nor.

	\noindent
	\Large {\bfseries \textcolor{red}{V/.}} \hspace{0.5cm} Demos gracias al Se\~nor, nuestro Dios.\\
	\noindent
	\Large {\bfseries \textcolor{red}{R/.}} \hspace{0.5cm} Es justo y necesario.

	\lettrine[lines=2]{\bfseries \textcolor{red}{E}}{}\Large n verdad es justo y necesario,\\
	es nuestro deber y salvación\\
	darte gracias siempre y en todo lugar,\\
	Señor, Padre santo, Dios todopoderoso y eterno,\\
	por Cristo, Señor nuestro.

	\noindent
	El cual, siendo inocente,\\
	se dignó padecer por los impíos,\\
	y ser condenado injustamente\\
	en lugar de los malhechores.\\
	De esta forma, al morir,\\
	borró nuestros delitos,\\
	y, al resucitar,\\
	logró nuestra salvación.

	\noindent
	\lettrine[lines=2]{\bfseries \textcolor{red}{P}}{}\Large or eso, te alabamos con todos los ángeles,\\
	aclamándote llenos de alegría:

	\large{\textcolor{red}{Con las manos juntas y en uni\'on del pueblo cantan o recitan en voz alta:}}

	\noindent
	\Large {Santo, Santo, Santo...}

	\Large {\bfseries \textcolor{red}{CONSAGRACI\'ON}}\\
	\large{\textcolor{red}{El celebrante, con las manos extendidas, dice:}}

	\lettrine[lines=2]{\bfseries \textcolor{red}{S}}{}\Large anto eres en verdad, Se\~nor,\\
	fuente de toda santidad.

	\newpage

	\large{\textcolor{red}{Junta las manos y, manteni\'endolas extendidas sobre las ofrendas, dice:}}

	\noindent
	\Large Santifiques estos dones\\
	con la efusi\'on de tu Esp\'iritu,

	\large{\textcolor{red}{Junta las manos y traza el signo de la cruz sobre el pan y el c\'aliz conjuntamente, diciendo:}}

	\noindent
	\Large de manera que sean para nosotros\\
	Cuerpo y \Huge{\textcolor{red}{\ding{64}}} \Large Sangre\\
	de Jesucristo, nuestro Se\~nor.

	\large{\textcolor{red}{Junta las manos.}}

	\lettrine[lines=2]{\bfseries \textcolor{red}{E}}{}\Large l cual,\\
	cuando iba a ser entregado a su Pasi\'on,\\
	voluntariamente aceptada,

	\large{\textcolor{red}{Toma el pan y, sosteni\'endolo un poco elevado sobre el altar, prosigue:}}

	\noindent
	\Large tom\'o pan,\\
	y elevando los ojos T\'i, Padre, Padre,\\
	pronuncio la bendici\'on, \\
	lo parti\'o y lo dio a sus disc\'ipulos, diciendo:

	\large{\textcolor{red}{Se inclina un poco.}}

	\noindent
	\LARGE{\bfseries{Tomad y comed todos de \'el,\\
	porque esto es mi Cuerpo,\\
	que ser\'a entregado por vosotros.}}

	\newpage

	\large{\textcolor{red}{Muestra el pan consagrado al pueblo, lo deposita luego sobre la patena y lo adora haciendo genuflexi\'on. Despu\'es prosigue:}}

	\noindent
	\Large Del mismo modo, acabada la cena,

	\large{\textcolor{red}{Toma el c\'aliz y, sosteni\'endolo un poco elevado sobre el altar, prosigue:}}

	\noindent
	\Large tom\'o el c\'aliz, lleno del fruto de la vid,\\
	y elevando los ojos a T\'i, Padre, Padre,\\
	pronuncio la bendici\'on,\\
	lo pas\'o a sus disc\'ipulos, diciendo:

	\large{\textcolor{red}{Se inclina un poco.}}

	\noindent
	\LARGE{\bfseries{Tomad y bebed todos de \'el,\\
	porque este es el c\'aliz de mi Sangre,\\
	Sangre de la alianza nueva y eterna,\\
	que ser\'a derramada por vosotros\\
	y por todos los hombres\\
	para el perd\'on de los pecados.\\
	Haced esto como mi memorial.}}

	\large{\textcolor{red}{Muestra el c\'aliz al pueblo, lo deposita luego sobre el corporal y lo adora haciendo genuflexi\'on.}}

	\noindent
	\Large {\bfseries \textcolor{red}{V/.}} \hspace{0.5cm} \'Este es el Sacramento de nuestra F\'e.

	\noindent
	\Large {\bfseries \textcolor{red}{R/.}} \hspace{0.5cm} Anunciamos tu muerte, Se\~nor,\\
	. \hspace{1.5cm} proclamamos tu resurrecci\'on.\\
	. \hspace{1.5cm} !`Maran-ath\'a! !`Maran-ath\'a! !`Maran-ath\'a!\\
	. \hspace{1.5cm} !`Ven, Se\~nor Jes\'us!

	\newpage

	\large{\textcolor{red}{Despu\'es el celebrante, con las manos extendidas, dice:}}

	\Large {\bfseries \textcolor{red}{ANAMNESIS}}

	\lettrine[lines=2]{\bfseries \textcolor{red}{A}}{}\Large s\'i, pues, Padre, al celebrar ahora\\
	el memorial de la muerte y resurrecci\'on de tu Hijo,\\
	te ofrecemos el pan de vida\\
	y el c\'aliz de salvaci\'on,\\
	y te damos gracias\\
	porque nos haces dignos de servirte en tu presencia.

	\lettrine[lines=2]{\bfseries \textcolor{red}{T}}{}\Large e pedimos humildemente\\
	que el Esp\'iritu Santo congregue en la unidad\\
	a cuantos participamos\\
	del Cuerpo y Sangre de Cristo.

	\noindent
	Acu\'erdate, Se\~nor, de tu Iglesia extendida por toda la tierra; \\
	y con el Papa {\bfseries \textcolor{red}{N.}},\\
	con nuestro Obispo {\bfseries \textcolor{red}{N.}},\\
	y todos los pastores que cuidan de tu pueblo,\\
	llevala a su perfecci\'on por la caridad.

	\lettrine[lines=2]{\bfseries \textcolor{red}{A}}{}\Large cu\'erdate tambi\'en de nuestros hermanos\\
	que se durmieron en la esperanza de la resurrecc\'ion,\\
	y de todos los que han muerto\\
	en tu misericordia;\\
	adm\'itelos a contemplar la luz de tu rostro.

	\noindent
	Ten misericordia de todos nosotros,\\
	y as\'i, con Mar\'a, la Virgen Madre de Dios,\\
	San Jos\'e su Fiel Esposo, \\
	los ap\'ostoles\\
	y cuantos vivieron en tu amistad\\
	a trav\'es de los tiempos,\\
	merezcamos, por tu Hijo Jesucristo,\\
	compartir la vida eterna\\
	y cantar tus alabanzas.

	\large{\textcolor{red}{Junta las manos. Toma la patena con el pan consagrado y el c\'aliz y, sosteni\'endolos elevados, dice:}}

	\lettrine[lines=2]{\bfseries \textcolor{red}{P}}{}\Large or Cristo, con \'el y en \'el,\\
	a ti, Dios Padre omnipotente,\\
	en la unidad del Esp\'iritu Santo,\\
	todo honor y toda gloria\\
	por los siglos de los siglos.

	\Large \hspace{-0.9cm} {\bfseries \textcolor{red}{R/.}} \hspace{0.5cm} Am\'en.

	\begin{center}
		\Huge {\bfseries Rito de la Comuni\'on}
	\end{center}

	\lettrine[lines=2]{\bfseries \textcolor{red}{F}}{}\Large ieles a la recomendaci\'on del Salvador\\
	y siguiendo su divina ense\~nanza,\\
	nos atrevemos a decir:

	\large{\textcolor{red}{Extiende las manos y, junto con el ministro, continua:}}

	\noindent
	\Large Padre nuestro...

	\large{\textcolor{red}{Solo el celebrante, con las manos extendidas, prosigue diciendo:}}

	\noindent
	\Large L\'ibramos Se\~nor de todos lo males,\\
	y conc\'edenos la paz en nuestros d\'ias,\\
	para que, ayudados por tu misericordia,\\
	vivamos siempre libres de pecado\\
	y protegidos de toda perturbaci\'on,\\
	mientras esperamos la gloriosa venida\\
	de nuestro Salvador Jesucristo.

	\large{\textcolor{red}{Junta las manos. El pueblo concluye la oraci\'on aclamando:}}

	\noindent
	\Large {\bfseries \textcolor{red}{R/.}} \hspace{0.5cm} Tuyo es el reino, tuyo es el poder y la gloria, por siempre, Se\~nor.

	\large{\textcolor{red}{Despu\'es el sacerdote, con las manos extendidas, dice en voz alta:}}

	\noindent
	\Large Se\~nor Jesucristo, que dijiste a tus Ap\'ostoles:\\
	<<La paz les dejo, mi paz les doy>>,\\
	no tengas en cuenta nuestros pecados,\\
	sino la fe de tu Iglesia,\\
	y conforme a tu palabra,\\
	conc\'edele la paz y la unidad.

	\large{\textcolor{red}{Junta las manos.}}

	\Large T\'u que vives y reinas por los siglos de los siglos.

	\noindent
	\Large {\bfseries \textcolor{red}{R/.}} \hspace{0.5cm} Am\'en.

	\begin{center}
	\Huge {\bfseries Rito de la Paz}
	\end{center}

	\large {\textcolor{red}{El celebrante, extendiendo y juntando las manos, dice:}}

	\noindent
	\Large {\bfseries \textcolor{red}{V/.}} \hspace{0.5cm} La paz del Se\~nor est\'e siempre con ustedes.\\
	\Large {\bfseries \textcolor{red}{R/.}} \hspace{0.5cm} Y con tu esp\'iritu.

	\noindent
	\Large {\bfseries \textcolor{red}{V/.}} \hspace{0.5cm} Dense fraternalmente la paz.

	\newpage

	\large{\textcolor{red}{Mientras el celebrante tama el pan consagrado, lo parte sobre la patena y pone una part\'icula dentro del caliz.}}

	\Large Cordero de Dios..

	\large{\textcolor{red}{El celebrante hace genuflexi\'on, toma el pan consagrado y, sosteni\'endolo un poco elevado, dice con voz clara:}}

	\noindent
	\Large \'Este es el Cordero de Dios,\\
	que quita el pecado del mundo.\\
	Dichosos los invitados a la cena del Se\~nor.

	\large{\textcolor{red}{Y, juntamente con el pueblo, a\~nade:}}

	\noindent
	\Large Se\~nor, no soy digno de que entres en mi casa,\\
	pero una palabra tuya bastar\'a para sanarme.

	\noindent
	\Large {\bfseries \textcolor{red}{R/.}} \hspace{0.5cm} Am\'en.

	\large{\textcolor{red}{Cuando el sacerdote comulga el Cuerpo de Cristo, comienza el canto de comuni\'on.}}

	\large{\textcolor{red}{Despu\'es el sacerdote puede ir a la sede. Si se juzga oportuno, se pueden guardar unos momentos de silencio o cantar un salmo o c\'antico de alabanza. Luego, de pie en la sede o en el altar, el sacerdote dice:}}

	\Large {\bfseries \textcolor{red}{Ant\'ifona de comuni\'on \hspace{1cm} Mt 26, 42}}\\
	\Large Padre mío, si este cáliz no puede pasar sin que yo lo beba, hágase tu voluntad.

	\Large {\bfseries \textcolor{red}{Oraci\'on despu\'es de la comuni\'on}}
	\lettrine[lines=2]{\bfseries \textcolor{red}{S}}{}\Large{ aciados con los dones santos,\\
	te pedimos, Señor, que,\\
	así como nos has hecho esperar\\
	lo que creemos por la muerte de tu Hijo,\\
	podamos alcanzar, por su resurrección,\\
	la plena posesión de lo que anhelamos.\\
	Por Jesucristo, nuestro Señor.}

	\noindent
	\Large {\bfseries \textcolor{red}{R/.}} \hspace{0.5cm} Am\'en.

	\Large {\bfseries \textcolor{red}{Oraci\'on despu\'es sobre el pueblo}}
	\lettrine[lines=2]{\bfseries \textcolor{red}{D}}{}\Large{ irige tu mirada, Señor,\\
	sobre esta familia tuya\\
	por la que nuestro Señor Jesucristo\\
	no dudó en entregarse a los verdugos\\
	y padecer el tormento de la cruz.\\
	Por Jesucristo, nuestro Señor.}

	\noindent
	\Large {\bfseries \textcolor{red}{R/.}} \hspace{0.5cm} Am\'en.

	\begin{center}
	\Huge {\bfseries Rito de Concluci\'on}
	\end{center}

	\large{\textcolor{red}{En este momento se hacen, si es necesario y con brevedad, los oportunos anuncios o advertencias al pueblo.}}

	\large{\textcolor{red}{Despu\'es tiene lugar la despedida. El sacerdote invita a la asamblea a inclinar la cabeza y bendice a todos diciendo:}}

	\noindent
	\Large {\bfseries \textcolor{red}{V/.}} \hspace{0.5cm} El Se\~nor este con vosotros.\\
	\noindent
	\Large {\bfseries \textcolor{red}{R/.}} \hspace{0.5cm} Y con tu Esp\'iritu.

	\noindent
	\Large {\bfseries \textcolor{red}{V/.}} \hspace{0.5cm} Dios, Padre de misericordia,\\
	. \hspace{1.5cm} que en la pasión de su Hijo os ha dado ejemplo de amor,\\
	. \hspace{1.5cm} os conceda, por vuestra entrega a Dios y a los hombres,\\
	. \hspace{1.5cm} la mejor de sus bendiciones.\\
	\Large {\bfseries \textcolor{red}{R/.}} \hspace{0.5cm} Am\'en.

	\newpage

	\noindent
	\Large {\bfseries \textcolor{red}{V/.}} \hspace{0.5cm} Y que gracias a la muerte temporal de Cristo,\\
	. \hspace{1.5cm} que alejó de vosotros la muerte eterna,\\
	. \hspace{1.5cm} obtengáis el don de una vida sin fin.\\
	\Large {\bfseries \textcolor{red}{R/.}} \hspace{0.5cm} Am\'en.

	\noindent
	\Large {\bfseries \textcolor{red}{V/.}} \hspace{0.5cm} Y así, imitando su ejemplo de humildad,\\
	. \hspace{1.5cm} participéis un día en su resurrección gloriosa.\\
	\Large {\bfseries \textcolor{red}{R/.}} \hspace{0.5cm} Am\'en.

	\noindent
	\Large {\bfseries \textcolor{red}{V/.}} \hspace{0.5cm} La bendici\'on de Dios todopoderoso,\\
	. \hspace{1.5cm} Padre, Hijo \Huge{\textcolor{red}{\ding{64}}} \Large y Esp\'iritu Santo,\\
	. \hspace{1.5cm} descienda sobre vosotros y os acompa\~ne siempre.\\
	\Large {\bfseries \textcolor{red}{R/.}} \hspace{0.5cm} Am\'en.

	\large{\textcolor{red}{Luego el celebrante, con las manos juntas, despide al pueblo:}}

	\noindent
	\Large {\bfseries \textcolor{red}{V/.}} \hspace{0.5cm} La Misa ha terminado. Pod\'eis ir en paz.\\
	\noindent
	\Large {\bfseries \textcolor{red}{R/.}} \hspace{0.5cm} Demos gracias a Dios.

	\newpage

	\pagecolor{red}\afterpage{\nopagecolor}
	\pagenumbering{gobble}

	\vspace*{\fill}
	\begin{center}
		\includegraphics{camino}
	\end{center}
	\vspace*{\fill}

% Termina el documento %%%%%%%%%%%%%%%%%%%%%%%%%%%%% F I N %%%%%%%%%%%%%%%%%%%%%%%%%%%%%%%%%%%%%%%%%%%%%%%%%%%%%%
\end{document}
